\section*{Chapter \ref{ch:b1:partial-equilibrium-welfare} --- Partial Equilibrium and Welfare}
% Numbers recomputed in code/bachelor-1/13-partial-equilibrium-welfare/exercises/solutions.py (tested).

\begin{solution}{exo:b1:partial-equilibrium-welfare:cs-ps}
$120 - 2p = p$: $p = 40$, $Q = 40$. Demand reaches zero at $p = 60$:
$\mathrm{CS} = \tfrac12 \times 40 \times 20 = 400$; $\mathrm{PS} = \tfrac12 \times 40 \times 40 = 800$.
\end{solution}

\begin{solution}{exo:b1:partial-equilibrium-welfare:read}
All items: $0.0 - 0.8 = -0.8$ points in July 2020 and $1.6 - (-0.7) = +2.3$ points in
January 2021. Food: $0.9 - 4.1 = -3.2$ points in July 2020.
\end{solution}

\begin{solution}{exo:b1:partial-equilibrium-welfare:share}
$\varepsilon_s/(\varepsilon_s + |\varepsilon_d|) = 1.5/2 = 0.75$: buyers bear three quarters.
\end{solution}

\begin{solution}{exo:b1:partial-equilibrium-welfare:cigarettes}
Buyers: their demand is inelastic, while sellers who can sell elsewhere
have elastic supply to any one market.
\end{solution}

\begin{solution}{exo:b1:partial-equilibrium-welfare:cs-integral}
$\int_1^2 100/p\,\dd p = 100\ln 2 = 69.31$.
\end{solution}

\begin{solution}{exo:b1:partial-equilibrium-welfare:exact}
Sellers keep $p$ and buyers pay $1.1p$: $100/(1.1p) = p$ gives $p = \sqrt{100/1.1} = 9.5346$,
$P = 10.4881$ and $Q = 9.5346$, against $p = Q = 10$ untaxed. The exact loss is
$\int_{9.5346}^{10}(100/x - x)\,\dd x = 100\ln(10/9.5346) - (100 - 90.909)/2 = 0.2201$. Both
elasticities are 1 in absolute value and $Q/p = 1$, so the formula gives
$\tfrac12\tau^2 \times \tfrac12$: $0.25$ with $\tau = 1$ and $0.2273$ with the actual wedge
$10.4881 - 9.5346 = 0.9535$. The formula is close; it errs because the curves
are not straight lines.
\end{solution}

\begin{solution}{exo:b1:partial-equilibrium-welfare:double}
To second order the loss is proportional to $\tau^2$: it is multiplied by
four. Revenue $\tau Q(\tau)$ less than doubles because the quantity falls as the
tax rises.
\end{solution}

\begin{solution}{exo:b1:partial-equilibrium-welfare:stability-ex}
$D = 20 + p$: \omterm{def:b1:models-and-data:equilibrium}{equilibrium} $p = 20$; $D' = 1 < S' = 2$, so $Z' = -1$ and
$\dot x = -0.5x$: \omterm{def:b1:partial-equilibrium-welfare:stable}{stable}. $D = -20 + 3p$: \omterm{def:b1:models-and-data:equilibrium}{equilibrium} $p = 20$; $D' = 3 > 2$, $Z' = 1$ and
$\dot x = 0.5x$: unstable, the price moves away from 20.
\end{solution}

\begin{solution}{exo:b1:partial-equilibrium-welfare:vat}
$1 - 1.16/1.19 = 2.52\%$ at the standard rate and $1 - 1.05/1.07 = 1.87\%$ at the
reduced rate. The all-items rate fell 0.8 points and the food rate 3.2
points from June to July 2020. The comparison is rough: the twelve-month
rate also moves with the change in prices of a year before and with every
other price change of the month, not all spending carries VAT (rents), and
pass-through was not full.
\end{solution}

\begin{solution}{exo:b1:partial-equilibrium-welfare:max}
$W(Q) = \int_0^Q(P^d - P^s)\,\dd x$ and $W'(Q) = P^d(Q) - P^s(Q)$, which is decreasing, zero at
the competitive quantity $Q^*$, positive before and negative after: $W$ rises
up to $Q^*$ and falls after it. Above $Q^*$ each extra unit costs more to
make than any buyer will pay, so producing it loses
$\int_{Q^*}^Q(P^s - P^d)\,\dd x > 0$.
\end{solution}

\begin{solution}{exo:b1:partial-equilibrium-welfare:substitute}
The tax raises the price of butter, so demand for margarine shifts out. With
upward-sloping supply, margarine's price rises: its producers gain surplus
(the area between the old and new price, left of the supply curve) and its
buyers lose about as much, since margarine is not taxed and its market stays
efficient: the change there is mostly a transfer. The dearer margarine in
turn raises the demand for butter, so butter buyers cut back less than they
would at the old margarine price. The right measure is the butter triangle
under a demand curve that lets the margarine price adjust, which is less
elastic; the triangle drawn with the margarine price held fixed
\emph{overstates} the loss.
\end{solution}

\begin{solution}{exo:b1:partial-equilibrium-welfare:refute}
Competitive incidence predicts pass-through between 0 and 1, symmetric for
a cut and a rise of the same size, larger where supply is more elastic
relative to demand. Across products: pass-through above one, or no
relation to the elasticities. Over time: prices that did not rise back
when the rates rose, or that moved by much more on the return than on the
cut.
\end{solution}

\begin{solution}{pb:b1:partial-equilibrium-welfare:sugar}
\textbf{1.} $120 - 10p = 10p$: $p = 6$, $Q = 60$.

\textbf{2.} $\varepsilon_d = -10 \times 6/60 = -1$ and $\varepsilon_s = 10 \times 6/60 = 1$.

\textbf{3.} $\dot p = k(120 - 20p)$; $Z'(p) = -20 < 0$, and
$p(t) = 6 + (p(0) - 6)e^{-20kt} \to 6$.

\textbf{4.} Demand reaches zero at $P = 12$: $\tfrac12 \times 60 \times 6 = 180$.

\textbf{5.} $\tfrac12 \times 60 \times 6 = 180$.

\textbf{6.} $W(Q) = \int_0^Q(12 - x/5)\,\dd x = 12Q - Q^2/10$: $W(60) = 360$, $W(50) = W(70) = 350$.

\textbf{7.} $120 - 10(p + 1.2) = 10p$: $p = 5.4$, $P = 6.6$, $Q = 54$.

\textbf{8.} $1/(1 + 1) = \tfrac12$: buyers pay $0.6$ more and sellers keep $0.6$ less, half each.

\textbf{9.} $1.2 \times 54 = 64.8$.

\textbf{10.} No: the wedge, the prices and the quantity are the same.

\textbf{11.} $\dd P/\dd\tau = 0.6/1.2 = 0.5$.

\textbf{12.} $p = P + 1.2$: $120 - 10P = 10(P + 1.2)$, so $P = 5.4$, $p = 6.6$ and $Q = 66$; it costs
$1.2 \times 66 = 79.2$.

\textbf{13.} $\tfrac12 \times 54 \times (12 - 6.6) = 145.8$.

\textbf{14.} $\tfrac12 \times 54 \times 5.4 = 145.8$.

\textbf{15.} $360 - 145.8 - 145.8 - 64.8 = 3.6$.

\textbf{16.} $\tfrac12 \times 1.2^2 \times (60/6) \times \tfrac{1 \times 1}{1 + 1} = 3.6$.

\textbf{17.} The curves are straight lines: the area lost is exactly a
triangle, and the second-order formula is exact.

\textbf{18.} $Q = 48$, revenue $2.4 \times 48 = 115.2$, loss $\tfrac12 \times 2.4 \times 12 = 14.4$: four times as
large, while revenue is less than twice as large.

\textbf{19.} $3.6/64.8 = 0.056$ and $14.4/115.2 = 0.125$.

\textbf{20.} $Q = 60 - 5\tau$, so $R(\tau) = \tau(60 - 5\tau)$.

\textbf{21.} $R'(\tau) = 60 - 10\tau = 0$ at $\tau = 6$: $P = 9$, $p = 3$, $Q = 30$ and $R = 180$.

\textbf{22.} $\mathrm{DWL} = 2.5 \times 36 = 90$: half a dalen per dalen raised.

\textbf{23.} $\mathrm{DWL}(\tau) = \tfrac12\tau \times 5\tau = 2.5\tau^2$, so $\mathrm{DWL}' = 5\tau$, while $R' = 60 - 10\tau$; the
loss per extra dalen raised is the ratio.

\textbf{24.} $6/48 = 0.125$ at $\tau = 1.2$ and $12/36 = 0.333$ at $\tau = 2.4$.

\textbf{25.} It grows without limit: near the top of the revenue curve a
higher tax raises almost no revenue but still destroys surplus.

\textbf{26.} At a $20\%$ rate the last dalen raised costs $0.125$ dalens of
\omterm{def:b1:partial-equilibrium-welfare:dwl}{deadweight loss}, more than twice the average loss per dalen, $0.056$.
\end{solution}
